\section{ANEXOS}

Los anexos deben publicarse en la página web del proyecto y ser los mismos especificados en la propuesta. Se relacionan los siguientes:

\begin{enumerate}[noitemsep]
    \item \textbf{Propuesta de trabajo de grado (VFP)}, versión 2.0 (13 de octubre de 2026).
    \item \textbf{Plan de gestión del proyecto (SPMP)}, IEEE 1058, con los modelos BPMN de ciclo de vida, requerimientos, riesgos, configuración y calidad, la estructura de desglose del trabajo y el cronograma.
    \item \textbf{Especificación de requisitos (SRS)}, versión 2.2, IEEE 830 e ISO/IEC/IEEE 29148: 60 requerimientos, casos de uso, modelo del dominio y plan de casos de prueba.
    \item \textbf{Documento de diseño (SDD)}, IEEE 1016: arquitectura, diseño de datos, de componentes y de interfaces, atributos de calidad y registros de decisiones (ADR).
    \item \textbf{Modelo C4:} diagramas de contexto, contenedores y componentes (\texttt{tesis.drawio} y sus exportaciones SVG), y diagramas de componentes y de clases en Mermaid (\texttt{c4-componentes-codigo.md}).
    \item \textbf{Documentación de la API:} especificación OpenAPI (\texttt{openapi.json}) y descripción de clases.
    \item \textbf{Código fuente:} repositorio con \texttt{simulation-service} (pipeline y API), \texttt{kepler} (visor), notebooks de exploración y datos procesados (\texttt{data/processed}).
    \item \textbf{Simulaciones precalculadas} de los cuatro presets (\texttt{simulation-service/precalculadas/}).
    \item \textbf{Informe de brechas de datos (F1.4)} y notebooks de las Fases 1 y 2.
    \item \textbf{Script de validación} \texttt{histograma\_salidas.py} y su figura.
    \item \textbf{Análisis de semaforización 3D} (\texttt{propuesta-semaforizacion-3d.md}).
    \item \textbf{Póster del proyecto.}
\end{enumerate}
